\documentclass[11pt]{article}
\usepackage[margin=1in]{geometry}
\usepackage{amsmath,amssymb,amsthm,mathtools}
\usepackage{graphicx}
\usepackage{xcolor}
\usepackage[hidelinks]{hyperref}
\usepackage{microtype}

\newtheorem{theorem}{Theorem}
\newtheorem{lemma}{Lemma}
\newcommand{\C}{\mathbb C}
\newcommand{\Sball}{\mathbb S}
\newcommand{\Dball}{\mathbb D}
\newcommand{\conv}{\operatorname{conv}}
\newcommand{\Rea}{\operatorname{Re}}

\title{The first nontrivial complex finite Grothendieck dimension is three}
\author{Automated mathematics run \texttt{fa\_banach\_001}\\
Lane 4, \texttt{agent\_lane\_04}, model GPT5.6}
\date{22 July 2026}

\begin{document}
\maketitle

\begin{center}
\fbox{\parbox{0.91\textwidth}{
\textbf{Status.} Candidate full solution, likely valid, pending human review.
The proof is exact; numerical optimization is used only as a sanity check.
}}
\end{center}

\section{Source question}

Yu, Duan, and Xu define the complex finite Grothendieck constant $K_d$ as
the least $K$ for which the unit ball $\Sball_d$ of $d\times d$ Schur
multipliers is contained in $K\conv(\Dball_d)$, where $\Dball_d$ consists
of rank-one Schur multipliers.  They recall that $K_2=1$ and ask for the
first integer $d$ such that $K_d>1$ \cite[p.~12]{YDX}:

\begin{quote}
``It seems, however, that the first integer $d$ for which $K_d>1$ is still
unknown.''
\end{quote}

\begin{figure}[h]
\centering
\includegraphics[width=0.93\textwidth]{figures/open_problem_crop.png}
\caption{The source question on PDF page 12 of arXiv:1201.1172.}
\end{figure}

The AQBC conjecture flagged by the deterministic scan is not the target of
this packet: it is disproved in the same source paper.  The first-dimension
question above is a separate, genuinely open statement in the source.

\section{Definitions and separation criterion}

We use inner products linear in the second variable.  The factorization
description of Schur multiplier norm gives
\[
 \Sball_d=\left\{S=(s_{ij}):s_{ij}=\langle \xi_i,\eta_j\rangle,
 \ \|\xi_i\|,\|\eta_j\|\leq1\right\},
\]
where the vectors may lie in any finite-dimensional Hilbert space.  The
rank-one set is
\[
 \Dball_d=\left\{(\overline{\alpha_i}\beta_j):
 |\alpha_i|,|\beta_j|\leq1\right\}.
\]
For a coefficient matrix $A=(a_{ij})$, define its vector and scalar support
values by
\begin{align*}
 V(A)&=\sup_{\|\xi_i\|,\|\eta_j\|\leq1}
 \left|\sum_{i,j}a_{ij}\langle\xi_i,\eta_j\rangle\right|,\\
 C(A)&=\sup_{|\alpha_i|,|\beta_j|\leq1}
 \left|\sum_{i,j}a_{ij}\overline{\alpha_i}\beta_j\right|.
\end{align*}
Thus $V(A)$ and $C(A)$ are the support functions of $\Sball_d$ and
$\conv(\Dball_d)$, respectively.  Consequently,
\begin{equation}\label{eq:separation}
 V(A)>C(A)\quad\Longrightarrow\quad K_d>1.
\end{equation}
This is simply strict separation of the two compact convex sets.

\section{Proof intuition}

The witness is a small Hermitian matrix with entries among
$\{1,-1,i,-i\}$.  A carefully chosen rank-two Gram matrix makes its three
vector row sums have squared norms $8/3,6,6$, producing the exact vector
value $8\sqrt6/3$.

For scalar phases, optimizing the three row phases leaves only two column
phases.  Cauchy--Schwarz and one weighted square reduce their total value to
a Hermitian three-cycle quadratic form.  Fixing one phase collapses the
quadratic maximization to the elementary function
$a+\sqrt{30-10a}$.  Its upper bound $11/2$ is sharp for the quadratic
relaxation, but the corresponding phases cannot satisfy equality in the
weighted square.  This incompatibility makes the scalar bound strict.

\section{Full solution}

\begin{theorem}\label{thm:main}
For the complex finite-dimensional constants in \cite{YDX}, one has
$K_3>1$.  Hence the first integer $d$ for which $K_d>1$ is $d=3$.
\end{theorem}

\begin{proof}
The equality $K_1=1$ is immediate from the definitions, and the known
assertion $K_2=1$ is recalled in the source passage above.  It is therefore
enough to prove $K_3>1$.

Consider
\begin{equation}\label{eq:A}
A=\begin{pmatrix}
1&1&1\\
1&-1&-i\\
1&i&-1
\end{pmatrix}.
\end{equation}
We first lower-bound $V(A)$.  Let $\eta_1,\eta_2,\eta_3$ be unit vectors
with Gram matrix
\begin{equation}\label{eq:G}
G=\begin{pmatrix}
1&(-5-7i)/18&(-5+7i)/18\\
(-5+7i)/18&1&(7-15i)/18\\
(-5-7i)/18&(7+15i)/18&1
\end{pmatrix}.
\end{equation}
This matrix is positive semidefinite: its eigenvalues are
\[
0,\qquad \frac32-\frac{\sqrt{179}}{18},\qquad
\frac32+\frac{\sqrt{179}}{18}.
\]
It therefore is the Gram matrix of three unit vectors in $\C^2$.  For the
$i$th row of $A$, put $w_i=\sum_j a_{ij}\eta_j$.  Direct multiplication
using \eqref{eq:G} gives
\begin{equation}\label{eq:wnorms}
\|w_1\|^2=\frac83,\qquad \|w_2\|^2=6,\qquad \|w_3\|^2=6.
\end{equation}
Taking $\xi_i=w_i/\|w_i\|$ yields
\begin{equation}\label{eq:vectorvalue}
V(A)\geq\sum_i\|w_i\|
=\sqrt{\frac83}+2\sqrt6=\frac{8\sqrt6}{3}.
\end{equation}

We now prove that the scalar support value is strictly smaller.  Since a
linear functional on each unit disk is maximized on its boundary, and a
common phase may be normalized, write the three column phases as $1,u,v$
with $|u|=|v|=1$.  Optimizing the row phases gives
\[
C(A)=\max_{|u|=|v|=1}\bigl(r_1+r_2+r_3\bigr),
\]
where
\begin{align*}
r_1&=|1+u+v|,\\
r_2&=|1-u-iv|,\\
r_3&=|1+iu-v|.
\end{align*}
Set $r=r_1$ and $Q=r_2^2+r_3^2$.  First,
\[
r_1+r_2+r_3\leq r+\sqrt{2Q}.
\]
The nonnegative square
$\bigl(\sqrt3\,r-\sqrt{2Q/3}\bigr)^2$ gives
\begin{equation}\label{eq:weighted}
\bigl(r+\sqrt{2Q}\bigr)^2\leq4r^2+\frac83Q,
\end{equation}
with equality only when $\sqrt{2Q}=3r$.

A direct expansion shows that the right-hand side of
\eqref{eq:weighted} equals
\begin{equation}\label{eq:H}
28+\frac83\Rea\left((1+2i)u+(1-2i)v
 +(3+4i)\overline u v\right).
\end{equation}
Put $a=\Rea((1+2i)u)$.  For fixed $u$, the real part inside
\eqref{eq:H} is at most
\[
a+\left|(1-2i)+(3+4i)\overline u\right|
=a+\sqrt{30-10a}.
\]
Since $a\leq\sqrt5<11/2$, squaring the desired inequality is legitimate,
and
\[
\left(\frac{11}{2}-a\right)^2-(30-10a)
=\left(a-\frac12\right)^2\geq0.
\]
Therefore
\begin{equation}\label{eq:quadratic-bound}
4r^2+\frac83Q\leq28+\frac83\cdot\frac{11}{2}
=\frac{128}{3}=\left(\frac{8\sqrt6}{3}\right)^2.
\end{equation}

It remains essential to verify strictness, since the quadratic bound itself
is sharp.  Equality in \eqref{eq:quadratic-bound} forces
\begin{equation}\label{eq:eqcase}
a=\frac12,\qquad
v=\frac{\overline{(1-2i)+(3+4i)\overline u}}{5}.
\end{equation}
The two unit-modulus solutions of $\Rea((1+2i)u)=1/2$ are
\[
u=\frac{(1-2i)(1\pm i\sqrt{19})}{10}.
\]
Substitution into \eqref{eq:eqcase} gives, respectively,
\begin{equation}\label{eq:r-equality-values}
r^2=\frac{112-24\sqrt{19}}{25}
\quad\hbox{or}\quad
r^2=\frac{112+24\sqrt{19}}{25}.
\end{equation}
On the other hand, if equality also held in \eqref{eq:weighted}, then
$2Q=9r^2$; combining this with equality in
\eqref{eq:quadratic-bound} would give
\[
\frac{128}{3}=4r^2+\frac83Q=16r^2,
\qquad\text{hence}\qquad r^2=\frac83.
\]
Neither value in \eqref{eq:r-equality-values} equals $8/3$.  Thus the
chain of scalar inequalities can never attain $128/3$.  The phase torus is
compact, so
\[
C(A)<\frac{8\sqrt6}{3}\leq V(A).
\]
The separation criterion \eqref{eq:separation} proves $K_3>1$, completing
the proof.
\end{proof}

\section{Verification notes}

The script \texttt{code/verify\_witness.py} checks symbolically that:
\begin{itemize}
\item $G$ is Hermitian with the displayed eigenvalues and rank two;
\item the three squared norms in \eqref{eq:wnorms} are exactly $8/3,6,6$;
\item the two equality cases give exactly the values in
  \eqref{eq:r-equality-values}, neither equal to $8/3$.
\end{itemize}
It also numerically maximizes the two-phase scalar expression as a sanity
check.  The numerical step is not used in the proof.

Potential reviewer failure modes are limited and explicit: the support
functional must use the same conjugation convention as the Schur
factorization, and strictness requires the final equality-case exclusion.

\section{Limitations and novelty check}

The theorem determines only the first dimension at which $K_d>1$.  It does
not compute the exact value of $K_3$, nor does it characterize all extreme
points of $\Sball_3$.

The bounded novelty search covered the run registry and solution indexes,
the local parsed arXiv corpus, and arXiv/web searches for the exact
first-integer phrase, complex $K_3$, and $3\times3$ extremal Schur
multipliers.  We also inspected Christensen's recent work
\cite{Christensen}, which proves necessary rank bounds for general extremal
Schur multipliers but does not supply this $K_3$ separation or answer the
source question.  No prior answer was located.  Novelty confidence is
moderate rather than definitive.

\section{Human-review recommendation}

\textbf{Recommend focused expert review.}  Check equations
\eqref{eq:wnorms}, \eqref{eq:H}, and the equality-case substitution
\eqref{eq:r-equality-values}.  If those three computations and the standard
support-function equivalence are confirmed, the source question is fully
answered.

\begin{thebibliography}{9}
\bibitem{YDX}
N. Yu, R. Duan, and Q. Xu,
\emph{Bounds on the distance between a unital quantum channel and the convex
hull of unitary channels, with applications to the asymptotic quantum
Birkhoff conjecture}, arXiv:1201.1172 (2012).

\bibitem{Pisier}
G. Pisier, \emph{Grothendieck's theorem, past and present},
Bull. Amer. Math. Soc. 49 (2012), 237--323.

\bibitem{Christensen}
E. Christensen, \emph{Extremal Schur Multipliers},
arXiv:2410.20112 (2024; latest inspected version dated March 2026).
\end{thebibliography}

\end{document}
